\chapter{The Vulkan Backend}
\label{ch:vulkan-backend}

Vulkan is CNA's second-most-mature backend by the project's own account, and the backend
whose history best illustrates a recurring theme in this Part: a class of bug where an
interface method accepts every parameter it is supposed to, but the implementation behind it
quietly discards most of them.

\section{Current test baseline}

At the time of writing, the full unit suite passes 4{,}371 of 4{,}373 tests (with two
hardware-dependent skips), and the Vulkan-specific \texttt{ctest} suite passes 126 of 127.
The single remaining failure, \texttt{Vulkan\_DepthBias}, is a narrow one: only the most
extreme tested magnitude of depth bias produces no observable effect
(\S\ref{sec:vulkan-depthbias}), a symptom shared with an equivalent finding on
\texttt{D3D9}. Historically this backend also carried five failing \cnaclass{BlendState}
tests and three segfaulting skinned-model content tests --- both fully resolved, with no
test exclusions needed to reach the current clean baseline.

\section{BlendState: from ``almost entirely fake'' to fully real}

The single most serious historical bug on this backend, in the project's own words, is that
its blend-state implementation was ``almost entirely fake'': one hardcoded blend equation was
applied regardless of what \cnaclass{BlendState} actually requested, confirmed failing across
five independent pixel tests before the fix. The fix replaced this with a genuine
per-\cnaclass{Blend} / \cnaclass{BlendFunction} mapping, applied consistently across all nine
of this backend's own pipeline-creation call sites --- meaning the fix had to be threaded
through every place a graphics pipeline object gets built, not just one central function.

\section{Stencil testing: the same ``discarded parameters'' bug shape}

\cnaclass{DepthStencilState} (Chapter~\ref{ch:state-objects}) had an almost identical
history. The backend's own state-application entry point accepted fifteen parameters but
genuinely stored only two of them (the depth-enable and depth-write flags) --- every
stencil-related parameter passed through the call was simply discarded, and the pipeline's
own stencil-test-enable flag was never set at all, on any pipeline. Five separate,
independently-run checks (the enable flag itself, the stencil masks, front-face operations,
two-sided stencil mode, and reference-stencil propagation) each separately confirmed this
failure mode. The fix added real per-pipeline depth-compare operations, full front- and
back-face \texttt{VkStencilOpState} configuration, and dynamic-state reference/mask values ---
and, as a direct side effect of the same fix, connected \texttt{ReferenceStencil}'s
independent-override path (Chapter~\ref{ch:state-objects}) on Vulkan specifically, something
still broken on EasyGL and BGFX. A second, smaller bug found in the same investigation: depth
format selection tried a stencil-\emph{less} format before it tried any stencil-capable
alternative.

\section{OcclusionQuery: from never-wired to a genuinely discriminating query}

Before its fix, occlusion queries on this backend were not merely inaccurate --- they were
never connected to anything at all. This backend defers 3D and 2D draws into a pending-batch
snapshot, recorded into real Vulkan commands only once per frame, well after
\texttt{Begin()} / \texttt{End()} had already returned synchronously; queries therefore always
reported zero, regardless of what was actually visible. The fix required resolving three
genuinely non-obvious design questions: how to tag an individual pending draw with the query
that should observe it; whether a single query may legitimately span multiple draw calls
(resolved: yes, as long as they share one render pass, tracked via contiguous-run detection,
with an explicit, documented limitation that a query spanning a render-pass \emph{boundary}
is not correctly summed --- a real capability gap against what a real device could do, stated
plainly rather than silently assumed); and how per-frame query-pool resets need to be
sequenced before any render pass begins. Verified in both directions --- a visible quad
producing a positive pixel count, an occluded quad behind a nearer opaque object producing
none --- plus a genuine multi-draw-span case (two non-overlapping half-quads inside one
\texttt{Begin} / \texttt{End} span, summing correctly), discriminating for real in this
project's own software Vulkan implementation.

\subsection{A worked example: using an OcclusionQuery correctly on this backend}

The real \cnaclass{OcclusionQuery} surface is small --- \texttt{Begin()}, \texttt{End()},
\texttt{getIsCompleteProperty()}, \texttt{getPixelCountProperty()} --- and the standard
visibility-culling pattern is exactly what the verification above tested directly:

\begin{lstlisting}
OcclusionQuery query(getGraphicsDeviceProperty());

query.Begin();
DrawBoundingProxyGeometry(occluderCandidate); // a cheap stand-in, not the real detailed mesh
query.End();

// One or more frames later, once rendering has actually reached the GPU:
if (query.getIsCompleteProperty() && query.getPixelCountProperty() > 0) {
    DrawFullDetailMesh(occluderCandidate); // visible -- worth the real draw cost
}
\end{lstlisting}

The one genuine gotcha this backend's own verification surfaced, worth restating as a rule
rather than only as a finding: a single query's \texttt{Begin()} / \texttt{End()} span may
legitimately cover \emph{multiple} draw calls (useful for querying a whole cluster of small
objects with one query instead of one per object), but only as long as every draw in that span
stays within one render pass. A query spanning a render-pass boundary --- switching render
targets, or anything else that closes and reopens a Vulkan render pass mid-span --- is not
correctly summed on this backend today. A conservative rule that avoids the gap entirely:
never let a render-target switch happen between an \cnaclass{OcclusionQuery}'s \texttt{Begin()}
and \texttt{End()}.

\section{Remaining limitation}
\label{sec:vulkan-depthbias}

Exactly one genuine, unresolved gap remains on this backend at the time of writing: an
isolated \cnaclass{RasterizerState.DepthBias} extreme-magnitude sub-case, where real hardware
bias produced no observable effect at any tested magnitude in this project's own sandboxed
environment --- suspected to be an environment or driver limitation shared with an equivalent
finding made independently on \texttt{D3D9} (Chapter~\ref{ch:direct3d-backends}), rather than
a Vulkan-specific defect.

\section{Two distinct render-target black-render bugs, both still open}
\label{sec:vulkan-rendertarget-black}

Chapter~\ref{ch:textures-rendertargets} previewed, in one sentence, that sampling a
\cnaclass{RenderTargetCube} face after rendering into it renders black on this backend ``for
two distinct, separately-tracked reasons.'' \texttt{docs/rendertarget-support.md} documents
both in enough detail to name precisely, and it is worth doing so here rather than leaving the
preview as the whole story.

The first (Task~875) is not specific to cube targets at all: \texttt{VulkanGraphicsBackend::%
Clear()} only records a global clear-color scalar and never registers the currently-bound
render target as actually ``used'' --- only a real draw call does that. The consequence is a
narrow but real trap for exactly the kind of code a cube-face render pass often looks like:

\begin{lstlisting}
device.SetRenderTarget(&probe, CubeMapFace::PositiveZ);
device.Clear(Color::CornflowerBlue); // records a clear color, but nothing else
device.SetRenderTarget(nullptr);     // no draw call happened in between -- Task 875's trigger
\end{lstlisting}

With no draw call between the \texttt{Clear} and the unbind, no Vulkan render pass is ever
recorded at all, and the target's underlying image stays permanently in
\texttt{VK\_IMAGE\_LAYOUT\_UNDEFINED} --- confirmed directly via a Vulkan validation error, not
inferred. Chapter~\ref{ch:textures-rendertargets}'s own \cnaclass{RenderTargetCube} worked
example sidesteps this specific trap by drawing real scene geometry into the probe (a genuine
draw call always follows the \texttt{Clear} there); a face that is deliberately left as flat
clear color with no further drawing is exactly the case this bug catches.

The second (Task~876) is the more interesting one, because it survives \emph{even after}
working around the first: rendering real content into all six \cnaclass{RenderTargetCube}
faces via an ordinary \cnaclass{SpriteBatch} draw (so Task~875 no longer applies), then
sampling that cube map back through \cnaclass{EnvironmentMapEffect} --- Chapter~\ref{ch:stock-effects}'s
own worked chrome-car-body example --- still renders black. The investigation's own account is
candid that this is not a simple, already-diagnosed bug: the sampling path itself
(\texttt{dynamic\_cast<IVulkanCubeSamplable*>} plus \texttt{GetVkCubeImageView()}) is described
as ``architecturally sound,'' unlike \textbf{BGFX}'s own cube-sampling bug covered in
Chapter~\ref{ch:bgfx-backend} (a genuine wrong-handle-cast, diagnosed precisely). Vulkan's root
cause instead remains genuinely unresolved between two named candidates. One possibility is that
the \cnaclass{SpriteBatch}-into-cube-face rendering itself is subtly wrong. The other is that the
descriptor-set cache behind environment-map sampling, \texttt{GetOrCreateEnvMapDescSet}, is
serving a stale or incorrect descriptor. The project's own record is explicit that this
ambiguity, not a lack of investigation, is why it remains open. A porting engineer relying on
real-time cube-map
reflections (dynamic environment mapping, a reflection probe updated every frame) should treat
this as a confirmed, currently-unworkaroundable gap on Vulkan specifically, not attempt to
debug their own call-site code against it first.

\section{A confirmed non-gap: device-feature gating}

Not every investigation into this backend found a bug. A dedicated check of whether optional
device features (\texttt{fillModeNonSolid}, needed for wireframe rendering, and
\texttt{samplerAnisotropy}) are properly negotiated before being requested found the gating
already correctly implemented --- both are queried via \texttt{vkGetPhysicalDeviceFeatures}
before being enabled, with a graceful fallback if the underlying device lacks either one.
This is included here deliberately: a book that only ever reports bugs found would give a
skewed picture of a backend whose overall health is, in fact, good.

\section{Notable feature-matrix specifics}

A cube-map render target's MSAA support is real but indirect: rather than a genuinely
multisampled cube attachment, it is implemented as an ordinary (non-cube) multisampled 2D
texture array paired with a separately cube-flagged resolve texture. Multiple render targets
are capped at four by shared, cross-backend code mirroring FNA's own
\texttt{MAX\_RENDERTARGET\_BINDINGS} constant --- notably \emph{not} because Vulkan's own
device capability is the limiting factor, since a typical device supports eight. Uploading
\texttt{Texture2D} data at a mip level above zero is a silent no-op on this backend (a shared
finding with Chapter~\ref{ch:state-objects}'s sampler-state coverage). Forwarding a real,
non-\texttt{Color} \cnaclass{SurfaceFormat} shares the identical blocked/architectural
limitation already covered for EasyGL in Chapter~\ref{ch:easygl-backend} --- both backends
are blocked by the same upstream \texttt{Texture::ValidateFormat} contract, not by anything
Vulkan-specific.

\section{Three tracked gaps, checked against live source rather than trusted from the doc}

This project's own \texttt{docs/rendertarget-support.md} --- the same document
\S\ref{sec:vulkan-rendertarget-black} already drew Tasks~875/876 from --- lists three further
tracked tasks for this backend as still open: real GPU mip generation for a render target
(Task~878), real render-target MSAA (Task~879), and any actual GPU effect from a custom
sub-region \cnaclass{Viewport} at all (Task~880, ``zero wiring on any backend''). Reading
\texttt{VulkanGraphicsBackend.cpp} directly, rather than trusting that document's own status
column, finds all three already resolved --- each superseded by a later task number
(\texttt{Task~903} / \texttt{907}) whose own comments the doc was never updated to reflect, one
more real instance of the recurring ``a status that looked settled turned out to be stale''
pattern this book has run into more than once. \texttt{VulkanRenderTargetBackend::%
MaybeGenerateMips()} runs a real, per-level \texttt{vkCmdBlitImage} cascade against a render
target's just-rendered (and, if applicable, just-MSAA-resolved) level-0 content immediately
after its render pass ends --- the direct Vulkan equivalent of EasyGL's own
\texttt{glGenerateMipmap}-on-unbind. Render-target MSAA is real for both
\cnaclass{RenderTarget2D} and \cnaclass{RenderTargetCube} (the cube case exactly as described
above). And \texttt{SetViewport()} genuinely reaches \texttt{vkCmdSetViewport} at
command-buffer-record time, mirroring \texttt{SetScissorRect}'s own identical storage-then-apply
shape --- with one honestly-scoped limitation worth stating precisely rather than rounding up to
``fixed'': a custom sub-region \cnaclass{Viewport} is only honored for the backbuffer pass. A
render-target pass stays hardcoded to that target's own full size, because this backend's
deferred, potentially-multiple-render-targets-per-frame recording model has no way to recover
\emph{which} \cnaclass{Viewport} was actually active at the moment each render target's own draws
were originally issued, by the time a frame's commands actually get recorded.

\section{A documented, not a hidden, coverage gap}

One honest gap in \emph{testing}, not rendering: at the time of writing there is no dedicated
pixel test for \cnaclass{SpriteBatch}'s sort-mode/rotation/scale/crop/flip behavior, for
\cnaclass{SpriteFont} glyph placement, or for a multi-mesh \cnaclass{Model} hierarchy on this
backend specifically --- confirmed by directly searching the test suite for matching test
names and finding none, despite basic \cnaclass{SpriteBatch}/2D-demo smoke tests passing.
This does not mean the underlying shared C\texttt{++} code is wrong (it is the same code
every other backend runs), only that this specific backend has not had these specific
features independently re-verified against its own rendering path.
